\documentclass[a4paper,10pt]{article}
\usepackage[czech]{babel}
\usepackage[T1]{fontenc}
\usepackage[utf8]{inputenc}
\usepackage{amsmath,amssymb}
\usepackage[margin=18mm]{geometry}
\setlength{\parindent}{0pt}
\begin{document}
\section*{Geometrická posloupnost A/B -- řešení a bodování}
Podklad pro učitele k testu
\texttt{01\_test\_GeometrickaPosloupnost\_AB.tex}.

Zařazení: 4. ročník SPŠ; tematický plán
\texttt{4\_rocnik/tematicke\_plany/SPS.pdf}, s.~1, září, bod 1.3.
Zářijový celek posloupností má společnou dotaci 12 hodin.
Vzory: úlohy 1, 19, 36 a 56 z pracovního listu:\\
\texttt{4\_rocnik/pracovni\_listy/pracovniListGeoPosl.tex}.
Předpokládaná délka testu je 20 minut, bez kalkulačky.
Jde o pracovní časový odhad, nikoli údaj ověřený ve třídě.

\subsection*{Varianta A}
\begin{enumerate}
\item $\frac{15}{5}=3$ a $\frac{45}{15}=3$. Oba podíly jsou stejné,
čísla tedy tvoří tři po sobě jdoucí členy geometrické posloupnosti s $q=3$.
\item Ze vztahu $a_4=a_1q^3$ plyne
\[
a_1=\frac{24}{2^3}=3,\qquad a_7=a_4q^3=24\cdot8=192.
\]
\item $a_1=3$, $q=2\ne1$, tedy
\[
S_6=a_1\frac{q^6-1}{q-1}=3\frac{64-1}{2-1}=189.
\]
\item Každá vyhovující trojice musí splňovat
\[
(x+1)^2=(x-1)(x+5)
\iff x^2+2x+1=x^2+4x-5
\iff 2x=6 \iff x=3.
\]
Zkouška v původním zadání: pro $x=3$ dostaneme $2,4,8$;
$4/2=8/4=2$. Hodnota vyhovuje a jiná možnost neexistuje.
Množina řešení je $K=\{3\}$.
\end{enumerate}

\subsection*{Varianta B}
\begin{enumerate}
\item $\frac{21}{7}=3$ a $\frac{63}{21}=3$. Oba podíly jsou stejné,
čísla tedy tvoří tři po sobě jdoucí členy geometrické posloupnosti s $q=3$.
\item Ze vztahu $a_4=a_1q^3$ plyne
\[
a_1=\frac{40}{2^3}=5,\qquad a_7=a_4q^3=40\cdot8=320.
\]
\item $a_1=5$, $q=2\ne1$, tedy
\[
S_6=a_1\frac{q^6-1}{q-1}=5\frac{64-1}{2-1}=315.
\]
\item Každá vyhovující trojice musí splňovat
\[
x^2=(x-3)(x+6)
\iff x^2=x^2+3x-18
\iff 3x=18 \iff x=6.
\]
Zkouška v původním zadání: pro $x=6$ dostaneme $3,6,12$;
$6/3=12/6=2$. Hodnota vyhovuje a jiná možnost neexistuje.
Množina řešení je $K=\{6\}$.
\end{enumerate}

\subsection*{Orientační bodování a kontrola}
Bodování je nezávazný návrh pro učitele; kritéria lze upravit podle
odborného posouzení. Uznávejte i jiné matematicky správné postupy.
\begin{enumerate}
\item 3 body: první podíl (1), druhý podíl (1), závěr o geometrické
posloupnosti se správným kvocientem (1).
\item 3 body: správný vztah pro první člen (1), hodnota $a_1$ (1),
hodnota požadovaného členu $a_7$ (1).
\item 3 body: vzorec pro součet geometrické řady (1), správné dosazení
a úprava mocniny (1), výsledek (1).
\item 4 body: vztah pro tři po sobě jdoucí členy (1), správná úprava
a výpočet $x$ (1), kvocient (1), ověření všech tří členů v původním zadání (1).
\end{enumerate}
Celkem 13 bodů v každé variantě. Stejnou početní chybu netrestejte opakovaně,
pokud navazující postup odpovídá získané hodnotě.
Orientační časy úloh v obou variantách: 2, 4, 4 a 6 minut;
zbývající přibližně 4 minuty jsou určeny na kontrolu.

Obě varianty mají stejné matematické postupy, indexy členů i délku součtu.
V úlohách 4 vede nutná podmínka k jediné lineární rovnici s nenulovým
koeficientem; zkouška zároveň prokazuje postačitelnost.
Nedělíme výrazy s proměnnou, takže se neztrácejí případné nulové členy.

Nezávislé ověření přesnou racionální aritmetikou: v úlohách 1 byly
porovnány oba podíly; v úlohách 2 byly členy vygenerovány postupným
násobením dvěma: A: $3,6,12,24,48,96,192$,
B: $5,10,20,40,80,160,320$.
V úlohách 3 dává přímé sčítání
$3+6+12+24+48+96=189$ a $5+10+20+40+80+160=315$.
V úlohách 4 byly zkontrolovány koeficienty rozdílu obou stran
($-2x+6$, resp. $-3x+18$) i podíly členů po dosazení.
\end{document}
